\documentclass[11pt]{article}
\usepackage[a4paper,margin=2.7cm]{geometry}
\usepackage{amsmath,amssymb}
\usepackage{graphicx}
\usepackage{booktabs}
\usepackage[hidelinks]{hyperref}
\usepackage{longtable}
\usepackage{float}
\usepackage{array}
\graphicspath{{figs/}}
\newcommand{\astar}{\alpha^{*}}
\renewcommand{\thesection}{S\arabic{section}}
\renewcommand{\thefigure}{S\arabic{figure}}
\renewcommand{\thetable}{S\arabic{table}}
\renewcommand{\theequation}{S\arabic{equation}}
\title{Supplementary Material for\\ ``Message Capacity and Claim Wording Set the Transition Points of Collective Truth-Finding in Language-Model Networks''}
\author{Makoto Fukushima\\[3pt]
{\normalsize Research Division, Honda Research Institute Japan Co., Ltd., Wako, Saitama, Japan}\\[2pt]
{\small makoto.fukushima@jp.honda-ri.com}}
\date{}
\begin{document}
\maketitle
\noindent Section numbers, figures, and tables in this document are prefixed with S. References to Sections I--VI, Figures 1--7, and Appendices A--D point to the main article. The glossary (Appendix B), the single-agent coefficient tables and the reduced-map validation (Appendix C), and the fidelity check (Appendix D) are in the main article. In this document and its tables, ``threshold'' denotes the transition point in the initial fraction (the unstable fixed point), which the main article calls the transition point.

\section{Reanalysis of the public dataset}

\subsection{Endogeneity of the initial fraction}

Variance decomposition of $x_0$ and within/between-question outcome
correlations (random-graph family, $\delta = 0.1$), computed over the 6{,}400
objective communities of El et al.:

\begin{center}
\begin{tabular}{lcccc}
\toprule
model & between-question share of $\mathrm{var}(x_0)$ &
$\mathrm{corr}(x_0,\text{outcome})$ & between & \textbf{within} \\
\midrule
Gemma-3n     & 0.995 & 0.802 & 0.852 & \textbf{0.041} \\
Llama-3-8B   & 0.936 & 0.659 & 0.721 & \textbf{0.162} \\
Qwen3.5-9B   & 0.861 & 0.277 & 0.351 & \textbf{0.132} \\
GPT-4o-mini  & 0.503 & 0.298 & 0.321 & \textbf{0.273} \\
\bottomrule
\end{tabular}
\end{center}

Question-clustered versus i.i.d.\ bootstrap intervals for the fitted
$q(x_0)$ crossing (random family):

\begin{center}
\begin{tabular}{lccccc}
\toprule
model & crossing & i.i.d.\ CI & width & question-clustered CI & ratio \\
\midrule
GPT-4o-mini & 0.494 & $[0.477, 0.513]$ & 0.036 & $[0.404, 0.583]$ & 4.9$\times$ \\
Gemma-3n    & 0.285 & $[0.251, 0.315]$ & 0.065 & $[0.163, 0.419]$ & 3.9$\times$ \\
Qwen3.5-9B  & 0.142 & $[0.059, 0.206]$ & 0.147 & $[-1.049, 0.367]$ & 9.6$\times$ \\
Llama-3-8B  & 0.459 & $[0.441, 0.480]$ & 0.039 & $[0.370, 0.564]$ & 5.0$\times$ \\
\bottomrule
\end{tabular}
\end{center}

Two illustrations of the endogeneity: communities at $x_0 \approx 0.40$ on
different questions run to opposite consensuses with near certainty, while
communities at $x_0 = 0.15$ on an easy question converge correctly in 32 of
32 episodes. These are limitations of the identification strategy, not of
the dataset's quality; the same logs support the message-level measurement
of the next subsection. The degenerate degree distributions noted in the
main article concentrate the random-family mass on $k = 2$--$4$, and the
lattice family has fixed degree.

\subsection{Message-level coefficients and the truth asymmetry}

Question-clustered 95\% intervals for the per-message truth asymmetry $w_T$
(43{,}500 observations per model):

\begin{center}
\begin{tabular}{lcccc}
\toprule
model & $w_T$ & cluster 95\% CI & Ising-coupling estimate & contains 1 \\
\midrule
GPT-4o-mini & 1.055 & $[0.996, 1.114]$ & 1.16 & yes (boundary) \\
Gemma-3n    & 1.179 & $[0.731, 1.628]$ & 1.71 & yes \\
Qwen3.5-9B  & 0.900 & $[0.759, 1.041]$ & 1.33 & yes \\
Llama-3-8B  & 1.144 & $[0.972, 1.316]$ & 1.38 & yes \\
\bottomrule
\end{tabular}
\end{center}

The label attenuation $\lambda_{\mathrm{lab}}$ (per-message weight of
``disagree''-labeled edges relative to ``agree''-labeled ones) is
$0.080 \pm 0.044$ (GPT-4o-mini), $0.555 \pm 0.078$ (Gemma), $0.294 \pm 0.051$
(Qwen), $0.288 \pm 0.041$ (Llama); $\theta$ is nonzero in all four models
($t \ge 3.4$, question-clustered).

\section{Main collective experiment}

\subsection{Distribution of final fractions}

The outcome labels of the main article ($x_8 \ge 0.6$ correct-side consensus,
$x_8 \le 0.4$ wrong-side) are operational thresholds. The table gives, per
campaign, the share of episodes that were decided and the share of decided
episodes whose final fraction lies beyond $|x_8 - 0.5| > 0.4$, with the
median end state on each side. Pooled over all 2{,}620 episodes, 66\% end in
$|x_8 - 0.5| \in [0.45, 0.5]$ (unanimity included).
\begin{center}
\small
\resizebox{\textwidth}{!}{\begin{tabular}{lrrrrr}
\toprule
campaign & episodes & decided ($|x_8-0.5|\ge 0.1$) & of which $|x_8-0.5|>0.4$ & median $x_8$, correct side & median $x_8$, wrong side \\
\midrule
main arm A ($\rho=0$) & 1008 & 0.895 & 0.634 & 0.922 & 0.031 \\
main arm A, rewired ($\rho=1$) & 96 & 0.969 & 0.667 & 1.000 & 0.031 \\
discrimination arm B & 300 & 0.917 & 0.633 & 0.875 & 0.062 \\
true claims (A-1) & 288 & 0.976 & 0.758 & 1.000 & --- \\
negated rewrites (A-2) & 192 & 1.000 & 0.995 & 1.000 & --- \\
qwen3:8b, eight claims & 256 & 0.996 & 0.992 & 1.000 & 0.000 \\
qwen3:8b, four claims & 480 & 0.996 & 0.967 & 1.000 & 0.000 \\
\bottomrule
\end{tabular}
}
\end{center}

\subsection{Outcome family (arm A, generated graphs)}
\begin{center}
\small
\begin{tabular}{rrrrccccc}
\toprule
$\alpha$ & $32x_0$ & $n$ & $q$ & Wilson 95\% & undec. & $q_{\delta=.05}$ & $q_{\delta=.2}$ & $q_{\text{excl}}$ \\
\midrule
0.2 & 4 & 16 & 0.000 & [0.000, 0.194] & 0.000 & 0.000 & 0.000 & 0.000 \\
0.2 & 8 & 40 & 0.025 & [0.004, 0.129] & 0.000 & 0.025 & 0.000 & 0.025 \\
0.2 & 12 & 16 & 0.000 & [0.000, 0.194] & 0.000 & 0.000 & 0.000 & 0.000 \\
0.2 & 16 & 16 & 0.000 & [0.000, 0.194] & 0.000 & 0.000 & 0.000 & 0.000 \\
0.2 & 20 & 40 & 0.100 & [0.040, 0.231] & 0.000 & 0.100 & 0.100 & 0.100 \\
0.2 & 24 & 40 & 0.275 & [0.161, 0.428] & 0.000 & 0.275 & 0.275 & 0.275 \\
0.3 & 4 & 40 & 0.000 & [-0.000, 0.088] & 0.025 & 0.000 & 0.000 & 0.000 \\
0.3 & 8 & 16 & 0.000 & [0.000, 0.194] & 0.062 & 0.062 & 0.000 & 0.000 \\
0.3 & 12 & 16 & 0.000 & [0.000, 0.194] & 0.000 & 0.000 & 0.000 & 0.000 \\
0.3 & 16 & 16 & 0.000 & [0.000, 0.194] & 0.125 & 0.000 & 0.000 & 0.000 \\
0.3 & 20 & 40 & 0.075 & [0.026, 0.199] & 0.200 & 0.100 & 0.050 & 0.094 \\
0.3 & 24 & 40 & 0.400 & [0.263, 0.554] & 0.075 & 0.425 & 0.400 & 0.432 \\
0.38 & 4 & 16 & 0.000 & [0.000, 0.194] & 0.000 & 0.000 & 0.000 & 0.000 \\
0.38 & 8 & 16 & 0.000 & [0.000, 0.194] & 0.062 & 0.000 & 0.000 & 0.000 \\
0.38 & 12 & 16 & 0.000 & [0.000, 0.194] & 0.125 & 0.000 & 0.000 & 0.000 \\
0.38 & 16 & 40 & 0.025 & [0.004, 0.129] & 0.075 & 0.025 & 0.000 & 0.027 \\
0.38 & 20 & 40 & 0.150 & [0.071, 0.291] & 0.050 & 0.200 & 0.150 & 0.158 \\
0.38 & 24 & 40 & 0.350 & [0.221, 0.505] & 0.125 & 0.350 & 0.325 & 0.400 \\
0.45 & 4 & 40 & 0.000 & [-0.000, 0.088] & 0.025 & 0.000 & 0.000 & 0.000 \\
0.45 & 8 & 40 & 0.000 & [-0.000, 0.088] & 0.150 & 0.000 & 0.000 & 0.000 \\
0.45 & 12 & 16 & 0.000 & [0.000, 0.194] & 0.125 & 0.000 & 0.000 & 0.000 \\
0.45 & 16 & 16 & 0.000 & [0.000, 0.194] & 0.125 & 0.000 & 0.000 & 0.000 \\
0.45 & 20 & 16 & 0.000 & [0.000, 0.194] & 0.312 & 0.062 & 0.000 & 0.000 \\
0.45 & 24 & 40 & 0.450 & [0.307, 0.602] & 0.075 & 0.475 & 0.450 & 0.486 \\
0.55 & 4 & 16 & 0.000 & [0.000, 0.194] & 0.062 & 0.000 & 0.000 & 0.000 \\
0.55 & 8 & 16 & 0.000 & [0.000, 0.194] & 0.125 & 0.000 & 0.000 & 0.000 \\
0.55 & 12 & 16 & 0.000 & [0.000, 0.194] & 0.125 & 0.000 & 0.000 & 0.000 \\
0.55 & 16 & 40 & 0.100 & [0.040, 0.231] & 0.175 & 0.125 & 0.050 & 0.121 \\
0.55 & 20 & 40 & 0.225 & [0.123, 0.375] & 0.175 & 0.225 & 0.150 & 0.273 \\
0.55 & 24 & 40 & 0.375 & [0.242, 0.530] & 0.075 & 0.400 & 0.300 & 0.405 \\
1 & 4 & 16 & 0.125 & [0.035, 0.360] & 0.188 & 0.188 & 0.125 & 0.154 \\
1 & 8 & 16 & 0.062 & [0.011, 0.283] & 0.125 & 0.125 & 0.000 & 0.071 \\
1 & 12 & 16 & 0.062 & [0.011, 0.283] & 0.312 & 0.125 & 0.000 & 0.091 \\
1 & 16 & 40 & 0.225 & [0.123, 0.375] & 0.275 & 0.300 & 0.150 & 0.310 \\
1 & 20 & 40 & 0.425 & [0.285, 0.578] & 0.125 & 0.450 & 0.300 & 0.486 \\
1 & 24 & 40 & 0.400 & [0.263, 0.554] & 0.275 & 0.525 & 0.325 & 0.552 \\
\bottomrule
\end{tabular}
\end{center}

\subsection{Empirical critical crowding level}
\begin{center}
\begin{tabular}{rrcc}
\toprule
$\alpha$ & $n$ (episodes, $x_0\le 8/32$) & wrong-consensus rate & claim-cluster 95\% CI \\
\midrule
0.2 & 56 & 0.982 & [0.946, 1.000] \\
0.3 & 56 & 0.964 & [0.893, 1.000] \\
0.38 & 32 & 0.969 & [0.906, 1.000] \\
0.45 & 80 & 0.912 & [0.775, 1.000] \\
0.55 & 32 & 0.906 & [0.719, 1.000] \\
1 & 32 & 0.750 & [0.438, 1.000] \\
\bottomrule
\end{tabular}
\end{center}
No crossing of the pre-registered 20\% threshold occurs in the window
($\astar_{\mathrm{emp}}$ undefined; 0 of 2{,}000 claim-clustered bootstrap
resamples cross).

\subsection{Discrimination arm (variant B)}
\begin{center}
\begin{tabular}{rrccc}
\toprule
$\alpha$ & $n$ & wrong rate ($x_0\le 10/32$) & 95\% CI & wrong basin retained \\
\midrule
0.6 & 40 & 0.975 & [0.925, 1.000] & True \\
0.9 & 40 & 1.000 & [1.000, 1.000] & True \\
1.3 & 40 & 0.975 & [0.925, 1.000] & True \\
\bottomrule
\end{tabular}
\end{center}

Observed low-$x_0$ ($x_0 \le 10/32$) mean final fractions against eight-round
surrogate rollouts (200 per episode, same graphs and seedings) of three rules
with claim-resolved fields: the normalized rule with the 17-claim identified
per-$k$ weights, and the per-claim reduced fits with divisive and with
constant coupling:
\begin{center}
\small
\resizebox{\textwidth}{!}{\begin{tabular}{lcccc}
\toprule
$\alpha$ & observed $\bar x_8$ & normalized (17-claim weights) & normalized (per-claim fit) & constant (per-claim fit) \\
\midrule
0.6 & 0.062 & 0.099 & 0.122 & 0.246 \\
0.9 & 0.066 & 0.133 & 0.164 & 0.323 \\
1.3 & 0.095 & 0.162 & 0.198 & 0.386 \\
\bottomrule
\end{tabular}
}
\end{center}
The normalized rules lie closer to the observations than the constant rule
but neither reproduces them exactly; claim 1569 converged wrong-side under
both rules' correct-side predictions (Section~IV-D of the main article).

\subsection{Rewiring control}
\begin{center}
\begin{tabular}{rrrrrr}
\toprule
$32x_0$ & $n_{\rho=0}$ & $q_{\rho=0}$ & $n_{\rho=1}$ & $q_{\rho=1}$ & $\Delta q$ \\
\midrule
4 & 40 & 0.000 & 16 & 0.000 & +0.0000 \\
8 & 16 & 0.000 & 16 & 0.000 & +0.0000 \\
12 & 16 & 0.000 & 16 & 0.000 & +0.0000 \\
16 & 16 & 0.000 & 16 & 0.000 & +0.0000 \\
20 & 40 & 0.075 & 16 & 0.062 & -0.0125 \\
24 & 40 & 0.400 & 16 & 0.312 & -0.0875 \\
\bottomrule
\end{tabular}
\end{center}

\subsection{Yoked one-step calibration}
Every realized transition replayed through the response function with the
17-claim identified per-$k$ weights and each claim's own field:
\begin{center}
\small
\begin{tabular}{lrrrrrr}
\toprule
arm & $\alpha$ & transitions & Brier & Brier (climatology) & ECE & residual at $l/k\in[0.4,0.6]$ \\
\midrule
A & 0.2 & 43,008 & 0.095 & 0.174 & 0.034 & -0.017 \\
A & 0.3 & 67,584 & 0.108 & 0.193 & 0.033 & +0.007 \\
A & 0.38 & 43,008 & 0.120 & 0.212 & 0.030 & +0.003 \\
A & 0.45 & 43,008 & 0.112 & 0.198 & 0.032 & -0.003 \\
A & 0.55 & 43,008 & 0.127 & 0.225 & 0.027 & +0.001 \\
A & 1 & 43,008 & 0.144 & 0.247 & 0.030 & -0.015 \\
B & 0.6 & 25,600 & 0.090 & 0.218 & 0.027 & -0.055 \\
B & 0.9 & 25,600 & 0.099 & 0.225 & 0.027 & -0.046 \\
B & 1.3 & 25,600 & 0.115 & 0.220 & 0.035 & -0.040 \\
\bottomrule
\end{tabular}

\end{center}
At contested compositions ($l/k \in [0.4, 0.6]$; last column) the residual
is within $\pm 0.02$ in arm A, whereas in arm B (own-state line present) a
systematic over-prediction of $0.04$--$0.06$ remains.

\section{Polarity discrimination experiment}

\subsection{SUPPORTS-claim collectives (A-1)}
\begin{center}
\begin{tabular}{llllll}
\toprule
claim & $\alpha$ & $32x_0=4$ & $32x_0=12$ & $32x_0=20$ & $32x_0=28$ \\
\midrule
993 & 0.3 & 0.820 & 1.000 & 1.000 & 1.000 \\
993 & 0.55 & 0.974 & 1.000 & 1.000 & 1.000 \\
1103 & 0.3 & 0.833 & 0.875 & 0.898 & 0.919 \\
1103 & 0.55 & 0.667 & 0.766 & 0.784 & 0.792 \\
1239 & 0.3 & 1.000 & 1.000 & 1.000 & 1.000 \\
1239 & 0.55 & 1.000 & 1.000 & 1.000 & 1.000 \\
\bottomrule
\end{tabular}
\end{center}

\subsection{Negated-rewrite collectives (A-2) versus originals}
\begin{center}
\small
\begin{tabular}{lllllll}
\toprule
rewrite & original & $32x_0=4$ & $32x_0=12$ & $32x_0=20$ & $32x_0=28$ & orig.\ mean $\bar x_8$ ($x_0\le 20/32$) \\
\midrule
90006 & 6 & 0.984 & 1.000 & 1.000 & 1.000 & 0.282 \\
90382 & 382 & 1.000 & 1.000 & 1.000 & 1.000 & 0.001 \\
91569 & 1569 & 1.000 & 1.000 & 1.000 & 1.000 & 0.223 \\
92070 & 2070 & 1.000 & 1.000 & 1.000 & 1.000 & 0.029 \\
\bottomrule
\end{tabular}
\end{center}

\subsection{Claim texts, rewrites, and checks}

Selected SUPPORTS claims (bare-prompt $p_q$ in parentheses): 993 (0.708)
``While there has been a mean rise of a little more than 3mm per year
worldwide since the 1990s, in the last decade, the NOAA Virginia Key tide
gauge just south of Miami Beach has measured a 9mm rise annually.''; 1103
(0.792) ``[ocean acidification was] First referenced in a peer-reviewed study
in Nature in 2003''; 1239 (0.792) ``Through decades of research, it has
become clear that human civilization, roughly 6,000 years old, developed
during an unusually stable period for global sea levels.''

Negated rewrites (model opposition-judgment rate; bare-prompt $p_q$ of the
rewrite, TRUE = correct side):

\begin{center}
\small
\begin{tabular}{p{5.6cm}p{5.6cm}cc}
\toprule
original & rewrite & opp.\ rate & $p_q$ \\
\midrule
The polar bear population has been growing. &
The polar bear population has \emph{not} been growing. & 0.70 & 0.344 \\
Sea level rise is decelerating. &
Sea level rise is \emph{not} decelerating. & 0.85 & 0.938 \\
The amount of energy used to construct solar and wind facilities is greater
than they produce in their working lives. &
\ldots is \emph{not} greater than they produce in their working lives. &
0.70 & 0.938 \\
When life is considered, ocean acidification is often found to be a
non-problem, or even a benefit. &
\ldots is often found to be a \emph{problem, and not a benefit}
(contrary; quantifier ``often'' retained). & 0.85 & 1.000 \\
\bottomrule
\end{tabular}
\end{center}

\subsection{Single-agent predictions for the polarity-experiment cells}

Per-claim divisive response functions (Appendix~A of the main article) give the predicted
phase of each (claim, $\alpha$) cell of the polarity experiment; all ten cells
are predicted monostable correct-side, and all ten converged correct-side.
In two cells (993 at $\alpha = 0.30$, 1103 at $\alpha = 0.55$) a minority of
episodes seeded at the lowest fraction ended wrong-side, so the fitted 50\%
crossing lies below the seeding grid; these are the two cells described in
the main article as matching to within the resolution of the grid.
\begin{center}
\small
\resizebox{\textwidth}{!}{\begin{tabular}{rrlll}
\toprule
claim & $\alpha$ & predicted (per-claim divisive) & observed & obs.\ $\bar x_8$ by $32x_0$ \\
\midrule
993 & 0.3 & mono-correct & correct from every seeded $x_0$ (crossing below grid) & 0.82 / 1.00 / 1.00 / 1.00 \\
993 & 0.55 & mono-correct & mono-correct & 0.97 / 1.00 / 1.00 / 1.00 \\
1103 & 0.3 & mono-correct & mono-correct & 0.83 / 0.88 / 0.90 / 0.92 \\
1103 & 0.55 & mono-correct & correct from every seeded $x_0$ (crossing below grid) & 0.67 / 0.77 / 0.78 / 0.79 \\
1239 & 0.3 & mono-correct & mono-correct & 1.00 / 1.00 / 1.00 / 1.00 \\
1239 & 0.55 & mono-correct & mono-correct & 1.00 / 1.00 / 1.00 / 1.00 \\
90006 & 0.3 & mono-correct & mono-correct & 0.98 / 1.00 / 1.00 / 1.00 \\
90382 & 0.3 & mono-correct & mono-correct & 1.00 / 1.00 / 1.00 / 1.00 \\
91569 & 0.3 & mono-correct & mono-correct & 1.00 / 1.00 / 1.00 / 1.00 \\
92070 & 0.3 & mono-correct & mono-correct & 1.00 / 1.00 / 1.00 / 1.00 \\
\bottomrule
\end{tabular}
}
\end{center}

\section{Cross-model campaigns}

\subsection{qwen3:8b: predicted fixed points versus observed collectives}

Fixed points of the mean-field map under each claim's divisive response
function (fitted from 336 single-shot queries per claim), with 95\%
single-agent-bootstrap intervals on the unstable point, against the observed
mean final fractions and the fitted 50\% crossing:
\begin{center}
\footnotesize
\resizebox{\textwidth}{!}{\begin{tabular}{rrllll}
\toprule
claim & $\alpha$ & predicted fixed points (s = stable, u = unstable) & pred.\ $x_c$ [95\% CI] & observed $\bar x_8$ ($32x_0 = 4/12/20/28$) & obs.\ crossing [95\% CI] \\
\midrule
6 & 0.3 & 0.0004 (s); 0.2087 (u); 1.0000 (s) & 0.21 [0.16, 0.28] & 0.200 / 1.000 / 1.000 / 1.000 & 0.19 [0.10, 0.25] \\
6 & 0.55 & 0.0010 (s); 0.1440 (u); 1.0000 (s) & 0.14 [0.09, 0.22] & 0.325 / 1.000 / 1.000 / 1.000 & 0.17 [0.08, 0.23] \\
6 & 1 & 1.0000 (s) & --- & 0.884 / 1.000 / 1.000 / 1.000 & below grid \\
1239 & 0.3 & 0.0000 (s); 0.6994 (u); 0.9998 (s) & 0.70 [0.64, 0.75] & 0.000 / 0.000 / 0.000 / 1.000 & 0.75 [0.72, 0.77] \\
1239 & 0.55 & 0.0000 (s); 0.7094 (u); 0.9997 (s) & 0.71 [0.64, 0.77] & 0.000 / 0.000 / 0.166 / 1.000 & 0.69 [0.63, 0.75] \\
1239 & 1 & 0.0000 (s); 0.7067 (u); 0.9995 (s) & 0.71 [0.64, 0.79] & 0.000 / 0.000 / 0.422 / 0.944 & 0.70 [0.61, 0.79] \\
1569 & 0.3 & 1.0000 (s) & --- & 1.000 / 1.000 / 1.000 / 1.000 & --- \\
1569 & 0.55 & 1.0000 (s) & --- & 1.000 / 1.000 / 1.000 / 1.000 & --- \\
1569 & 1 & 1.0000 (s) & --- & 1.000 / 1.000 / 1.000 / 1.000 & --- \\
91569 & 0.3 & 0.0195 (s); 0.3738 (u); 0.9940 (s) & 0.37 [0.22, 0.49] & 0.178 / 0.656 / 1.000 / 1.000 & 0.35 [0.28, 0.43] \\
91569 & 0.55 & 0.9958 (s) & --- & 0.953 / 1.000 / 1.000 / 1.000 & --- \\
91569 & 1 & 0.9973 (s) & --- & 1.000 / 1.000 / 1.000 / 1.000 & --- \\
\bottomrule
\end{tabular}
}
\end{center}

Per-cell outcome counts (correct/wrong/undecided of 10 episodes) for the same
campaign; intermediate mean final fractions in Fig.~6 of the main article are
splits between the two basins:
\begin{center}
\small
\begin{tabular}{rrcccc}
\toprule
claim & $\alpha$ & $32x_0=4$ & $32x_0=12$ & $32x_0=20$ & $32x_0=28$ \\
\midrule
6 & 0.3 & 2/8/0 & 10/0/0 & 10/0/0 & 10/0/0 \\
6 & 0.55 & 3/7/0 & 10/0/0 & 10/0/0 & 10/0/0 \\
6 & 1 & 9/1/0 & 10/0/0 & 10/0/0 & 10/0/0 \\
1239 & 0.3 & 0/10/0 & 0/10/0 & 0/10/0 & 10/0/0 \\
1239 & 0.55 & 0/10/0 & 0/10/0 & 2/8/0 & 10/0/0 \\
1239 & 1 & 0/10/0 & 0/10/0 & 3/7/0 & 9/0/1 \\
1569 & 0.3 & 10/0/0 & 10/0/0 & 10/0/0 & 10/0/0 \\
1569 & 0.55 & 10/0/0 & 10/0/0 & 10/0/0 & 10/0/0 \\
1569 & 1 & 10/0/0 & 10/0/0 & 10/0/0 & 10/0/0 \\
91569 & 0.3 & 0/9/1 & 6/4/0 & 10/0/0 & 10/0/0 \\
91569 & 0.55 & 10/0/0 & 10/0/0 & 10/0/0 & 10/0/0 \\
91569 & 1 & 10/0/0 & 10/0/0 & 10/0/0 & 10/0/0 \\
\bottomrule
\end{tabular}

\end{center}

qwen3:8b calibration of the 120 candidate claims found 4 SUPPORTS and 3
REFUTES claims with $|h_q| < 0.4$; a polarity-balanced set of 12 could not be
formed, so genuine truth asymmetry remains unidentified in this model as
well.

\subsection{llama3.3:70b single-agent probe}
\begin{center}
\begin{tabular}{rrrrrr}
\toprule
claim & $c_0$ & $\beta_T$ & $\beta_F$ & $n$ & $P(\text{TRUE})$ at $k=0$ \\
\midrule
6 & 1.987 & 0.636 & -0.257 & 132 & 0.000 \\
382 & 3.920 & 16.536 & -0.500 & 132 & 0.000 \\
1569 & 36.547 & 5.377 & -4.482 & 132 & 0.000 \\
2070 & 3.455 & 1.375 & -0.554 & 132 & 0.000 \\
90006 & 0.378 & 0.487 & -0.494 & 132 & 0.417 \\
90382 & 1.739 & 1.407 & -0.748 & 132 & 1.000 \\
91569 & 1.713 & 1.156 & -0.572 & 132 & 1.000 \\
92070 & 3.910 & 0.203 & 0.203 & 132 & 1.000 \\
\bottomrule
\end{tabular}
\end{center}
Polarity-coded regression across all eight claims (claim fixed effects):
$b(\text{TRUE-side message}) = +0.583$, $b(\text{FALSE-side}) = -0.563$,
ratio 1.03 with claim-clustered 95\% interval $[0.66, 1.64]$ (2{,}000
resamples of the eight claims; weak ridge $10^{-3}$ against separation). On
the same eight claims and message bank llama3.1:8b gives 1.94 $[1.39, 2.73]$
(over the 17 calibration claims: 1.44 $[1.25, 1.65]$); the 8b$-$70B
difference is 0.90 $[0.00, 1.72]$ and the ratio of ratios 1.87
$[1.00, 3.12]$, so the reduction is at the edge of resolution with eight
claims. Per-claim coefficients with $|c_0| > 5$ in the table above reflect
separation (accuracy 0 or 1 at $k = 0$) and are reported for completeness.
Empty-inbox fields (last column, observed rates over the empty-inbox
queries of each claim): the four false originals give
$P(\mathrm{TRUE}) = 0.000$ (every answer on the denying and correct side); of
the four true rewrites, three give 1.000 and the negated polar-bear rewrite
(90006) 0.42, the one near-neutral claim among the eight.
Per-$k$ fits (truth-coded, pooled across claims):
$\beta_T(4) = 2.41 \to \beta_T(8) = 1.35$. The positive wrong-side
coefficient at $k = 4$ ($+0.61$) --- wrong-side messages slightly
\emph{increasing} correct-side probability --- is recorded as an observation
only.

\section{Stage 2c: eight-claim calibration campaign and polarity-origin probe}

\subsection{Calibration campaign: all 16 cells}

Predicted classifications and thresholds come from each claim's divisive
response function fitted to its single-shot queries (Appendix~A of the main article); the eight
claims were selected from twenty candidates by the pre-registered rule.
Prediction intervals are single-agent bootstrap percentiles (200 resamples
of each claim's queries within design cells) and ``frac.\ bistable'' is the
fraction of resamples classified bistable. Observed crossings are 50\% points
of ridge-penalized logistic fits of episode outcomes against $x_0$, with
episode-bootstrap 95\% intervals;
$n_{\mathrm{c}}/n_{\mathrm{w}}/n_{\mathrm{u}}$ are correct/wrong/undecided
episode counts of 16 per cell.

\begin{table}[h]
\centering\footnotesize
\resizebox{\textwidth}{!}{\begin{tabular}{rrlllll}
\toprule
claim & $\alpha$ & predicted & pred.\ $x_c$ [95\% CI] (frac.\ bistable) & observed & obs.\ crossing [95\% CI] & $n_{\mathrm{c}}/n_{\mathrm{w}}/n_{\mathrm{u}}$ \\
\midrule
97 & 0.3 & mono-correct & --- & mono-correct & --- & 16/0/0 \\
97 & 0.55 & mono-correct & --- & mono-correct & --- & 16/0/0 \\
184 & 0.3 & mono-correct & --- & mono-correct & --- & 16/0/0 \\
184 & 0.55 & mono-correct & --- & mono-correct & --- & 16/0/0 \\
590 & 0.3 & bistable & 0.559 [0.497, 0.619] (1.00) & bistable & 0.560 [0.465, 0.653] & 7/9/0 \\
590 & 0.55 & bistable & 0.578 [0.516, 0.637] (1.00) & bistable & 0.622 [0.536, 0.701] & 6/10/0 \\
656 & 0.3 & bistable & 0.185 [0.125, 0.267] (1.00) & bistable & 0.248 [0.157, 0.341] & 12/4/0 \\
656 & 0.55 & bistable & 0.119 [0.072, 0.187] (1.00) & bistable & 0.120 [0.057, 0.200] & 14/2/0 \\
832 & 0.3 & bistable & 0.425 [0.364, 0.494] (1.00) & bistable & 0.498 [0.413, 0.589] & 8/8/0 \\
832 & 0.55 & bistable & 0.396 [0.331, 0.466] (1.00) & bistable & 0.436 [0.334, 0.531] & 9/7/0 \\
1013 & 0.3 & bistable & 0.407 [0.344, 0.470] (1.00) & bistable & 0.436 [0.335, 0.525] & 9/7/0 \\
1013 & 0.55 & bistable & 0.409 [0.346, 0.471] (1.00) & bistable & 0.373 [0.292, 0.460] & 10/6/0 \\
1146 & 0.3 & bistable & 0.094 [0.046, 0.157] (0.98) & mono-correct & --- & 16/0/0 \\
1146 & 0.55 & bistable & 0.034 [0.025, 0.101] (0.62) & bistable & 0.012 [-0.087, 0.169] & 15/1/0 \\
1151 & 0.3 & bistable & 0.271 [0.147, 0.395] (0.96) & bistable & 0.184 [0.085, 0.285] & 13/3/0 \\
1151 & 0.55 & bistable & 0.288 [0.178, 0.407] (0.95) & bistable & 0.311 [0.208, 0.404] & 11/4/1 \\
\bottomrule
\end{tabular}
}
\caption{All 16 cells of the eight-claim calibration campaign (qwen3:8b,
256 episodes, zero parse failures). Fifteen cells match in class; the
mismatch is claim 1146 at $\alpha = 0.30$, predicted bistable with threshold
0.094 just above the lowest seeded fraction (2/32), observed monostable
correct. All eleven predicted thresholds of the cells observed bistable lie
inside the observed intervals (mean absolute difference 0.04; Spearman
$r_s = 0.96$, claim-permutation $p = 0.004$).}
\end{table}

\subsection{Polarity-origin probe (D-3)}

Judgment-side single-shot measurements on three llama3.1-8b variants with the
fixed 8B message bank and a three-shot answer-format prefix applied
identically to all variants. The base model does not follow the forced-choice
format (parse rate 1.5\% in the bare pilot, 4.9\% under the prefix), so its row is
reported for completeness and carries no statistical weight.

\begin{table}[h]
\centering\small
\resizebox{\textwidth}{!}{%
\begin{tabular}{lllll}
\toprule
variant & tag (digest) & parsed & $|b_T|/|b_F|$ & 95\% CI \\
\midrule
instruct (control) & llama3.1:8b (46e0c10c039e) & 5{,}984/5{,}984 & 1.437 & [1.250, 1.651] \\
abliterated & mannix/llama3.1-8b-abliterated (f7d729db9832, Q4\_0) & 5{,}984/5{,}984 & 1.243 & [1.051, 1.471] \\
base & llama3.1:8b-text-q4\_K\_M (6f98b5a6e4b7) & 294/5{,}984 & (2.10) & not evaluable \\
\bottomrule
\end{tabular}}
\caption{Per-message assertion asymmetry (claim-fixed-effect logistic
regression; claim-clustered bootstrap intervals). The instruct control
reproduces the Stage-1 value 1.44 under the added format prefix. Bare
affirmation indices $S$ for the four original/negated pairs: instruct
0.92--1.52 (mean 1.22), abliterated 1.13--1.42 (mean 1.24); base cells parsed
2--6 of 64 and are not evaluable.}
\end{table}

\section{Residual and consistency analyses}

\subsection{Message statistics by side}
\begin{center}
\begin{tabular}{rlrrr}
\toprule
claim & side & messages & words/msg & hedges/msg \\
\midrule
6 & correct & 28258 & 60.7 & 0.058 \\
6 & wrong & 35998 & 59.4 & 0.419 \\
382 & correct & 16713 & 63.5 & 0.035 \\
382 & wrong & 47591 & 62.9 & 0.116 \\
1569 & correct & 28115 & 58.8 & 0.084 \\
1569 & wrong & 36093 & 59.6 & 0.314 \\
2070 & correct & 14105 & 64.5 & 0.081 \\
2070 & wrong & 50119 & 59.8 & 1.494 \\
\bottomrule
\end{tabular}
\end{center}

\subsection{Recorded-message persuasiveness probes}
\begin{center}
\begin{tabular}{rllcr}
\toprule
claim & side & message source & $P$(correct) & $n$ \\
\midrule
6 & correct & bank & 1.000 & 60 \\
6 & correct & episode & 1.000 & 60 \\
6 & wrong & bank & 0.100 & 60 \\
6 & wrong & episode & 0.167 & 60 \\
382 & correct & bank & 1.000 & 60 \\
382 & correct & episode & 1.000 & 60 \\
382 & wrong & bank & 0.017 & 60 \\
382 & wrong & episode & 0.000 & 60 \\
1569 & correct & bank & 1.000 & 60 \\
1569 & correct & episode & 1.000 & 60 \\
1569 & wrong & bank & 0.100 & 60 \\
1569 & wrong & episode & 0.050 & 60 \\
2070 & correct & bank & 1.000 & 60 \\
2070 & correct & episode & 1.000 & 60 \\
2070 & wrong & bank & 0.050 & 60 \\
2070 & wrong & episode & 0.167 & 60 \\
\bottomrule
\end{tabular}
\end{center}
Episode-generated and bank messages are indistinguishable in persuasive
force; no message-quality contribution to the high-$x_0$ residual is detected.

\subsection{Repeated exposure and finite time}

One-step regression with round interactions gives
$\beta_T \times (t-1) = +0.026$ (SE 0.006) and $\beta_F \times (t-1) =
-0.004$ (SE 0.003): repeated exposure does not attenuate message influence.
Finite-time check on high-$x_0$ episodes:
\begin{center}
\small
\begin{tabular}{rrrrrr}
\toprule
$\alpha$ & $32x_0$ & $n$ & mean $x_8-x_7$ & frac.\ interior & frac.\ still moving \\
\midrule
0.2 & 20 & 40 & -0.008 & 0.300 & 0.275 \\
0.2 & 24 & 40 & -0.009 & 0.225 & 0.150 \\
0.3 & 20 & 40 & -0.012 & 0.450 & 0.250 \\
0.3 & 24 & 40 & +0.001 & 0.275 & 0.325 \\
0.38 & 20 & 40 & +0.001 & 0.450 & 0.375 \\
0.38 & 24 & 40 & -0.028 & 0.375 & 0.350 \\
0.45 & 20 & 16 & -0.020 & 0.438 & 0.375 \\
0.45 & 24 & 40 & -0.008 & 0.275 & 0.375 \\
0.55 & 20 & 40 & -0.009 & 0.525 & 0.375 \\
0.55 & 24 & 40 & -0.002 & 0.450 & 0.225 \\
1 & 20 & 40 & -0.002 & 0.800 & 0.575 \\
1 & 24 & 40 & -0.027 & 0.725 & 0.550 \\
\bottomrule
\end{tabular}
\end{center}

\subsection{Label residual}

Correlation between an episode's share of TF label assignments and its final
state: $r = 0.194$ ($n = 1{,}008$ episodes) --- an order of magnitude below
the experimental effects ($\Delta \bar x_8 \approx 1$), but nonzero;
per-call randomization protects pooled estimates.

\subsection{Flip statistics and cascades}

Refractory regression: the coefficient of ``flipped at $t-1$'' on flipping at
$t$ is $+0.98$ (claim-cluster SE 0.13) --- flips concentrate in particular
agents (anti-refractory). Branching ratios:
\begin{center}
\begin{tabular}{rrr}
\toprule
$\alpha$ & flip events & $\sigma$ (mean downstream flips) \\
\midrule
0.2 & 9048 & 3.546 \\
0.3 & 10129 & 3.102 \\
0.38 & 11452 & 2.771 \\
0.45 & 10473 & 2.476 \\
0.55 & 12301 & 2.312 \\
1 & 14625 & 1.776 \\
\bottomrule
\end{tabular}
\end{center}
Cascade-size distribution in the asynchronous public dataset (GPT-4o-mini,
320 runs; cascades linked by graph adjacency within three events):
\begin{center}
\begin{tabular}{rr}
\toprule
cascade size & count \\
\midrule
1 & 9341 \\
2 & 668 \\
3 & 130 \\
4 & 28 \\
5 & 9 \\
6 & 1 \\
7 & 1 \\
\bottomrule
\end{tabular}
\end{center}
91.8\% of cascades have size 1; the tail decays geometrically (deeply
subcritical).

\section{Extended Methods}

This section reproduces in full the Methods subsections that Appendix~A of the main article
summarizes: models and serving, verbatim prompts, the cross-model
calibration campaign, the complete statistical procedures, and the
pre-registration statement.

\subsection{Models and serving}

All agents are served with ollama. The primary model is llama3.1:8b
(Q4\_K\_M quantization, temperature 0.7 throughout, seed recorded for every
call). The cross-model campaign uses qwen3:8b (tag 500a1f067a9f, thinking
disabled via \texttt{think:false}) and llama3.3:70b (Q4\_K\_M; judgments
only, message banks fixed to the 8B-generated banks). Single-shot campaigns
ran against a user-space ollama instance with \texttt{OLLAMA\_NUM\_PARALLEL=16}
(the system instance saturates at 0.62 generations/s; the parallel instance
reaches 3.5 generations/s); the 70B model does not load under 16-slot KV
allocation and was served by the serial system instance. The polarity-origin probe (Section~V of the main article) used two further variants, judgments
only: llama3.1:8b-text-q4\_K\_M (the base model) and an abliterated derivative
of the instruct model (mannix/llama3.1-8b-abliterated, Q4\_0), with a
three-shot answer-format prefix applied identically to all variants,
including the re-measured instruct control. Every call's prompt
hash, seed, raw response, and parse result are logged to JSONL.

\subsection{Prompts}

The judgment prompt (verbatim; variant A) is:
\begin{quote}\footnotesize\ttfamily
\# Your Expertise\\
\{persona\}\\[2pt]
\# Messages\\
You received the following messages from fellow panelists:\\
- \{message 1\} \ldots\\[2pt]
\# Problem\\
You are one member of an expert panel judging the following claim.\\
Claim: "\{claim\}"\\
Options:\\
A) \{TRUE or FALSE\}\\
B) \{FALSE or TRUE\}\\
\# Instruction\\
Give your direct answer on which option is correct.\\
Do NOT work through the problem, show any steps, compute, or\\
explain --- just commit to a snap judgement.\\
\#\# Response Format\\
Respond with a single character, A or B, and nothing else.
\end{quote}
Variant B inserts, before \texttt{\# Problem}: ``\texttt{\# Your Current
Answer / Your current answer is that the claim is \{TRUE|FALSE\}.}'' The A/B
$\leftrightarrow$ TRUE/FALSE assignment is randomized per call and
counterbalanced within cells; message presentation order is randomized and
recorded. The message-generation prompt asks the sender, given its persona,
the claim, and its current stance, to ``write a brief message to a fellow
panelist stating that you believe the claim is \{stance\} and giving the
single strongest reason from your reasoning,'' with a two-sentence response
format; stance is expressed in the message body (TRUE/FALSE), never as an
option letter, because option labels are randomized per receiver. Personas are
one-line expertise statements drawn from a fixed pool of 12. Prompt texts are
extracted verbatim from the released scripts.

\subsection{Cross-model calibration campaign}

Twenty further CLIMATE-FEVER claims (ten SUPPORTS, ten REFUTES, disjoint from
all earlier sets) were calibrated under the full scaffold and measured
single-shot with qwen3:8b (24 calls per candidate, 64 in the boundary band;
per-claim response-function fits from 6{,}720 further queries). The resulting
fixed-point classifications and thresholds for $\alpha \in \{0.30, 0.55\}$
were computed at the time of the campaign, and eight claims were selected by
a rule fixed in the pre-registered instructions (at least four predicted bistable with interior
thresholds, the remainder monostable, filled from the bistable pool where a
category was empty --- no wrong-side-dominant claim existed among the twenty
candidates). Collectives ran the full grid $\alpha \in \{0.30, 0.55\} \times
x_0 \in \{2, 6, \ldots, 30\}/32$ with two replicates (256 episodes, protocol
identical to the main cross-model campaign). Observed thresholds are defined
under Statistics (Appendix~A of the main article) and in the next subsection.

\subsection{Statistics: full procedures}

All confidence intervals on measured coefficients are claim-clustered pairs
bootstrap percentiles (1{,}000--2{,}000 resamples for the pooled
coefficients, 200 for the identified per-$k$ coefficients and for the
prediction-side intervals, 2{,}000 for the polarity ratios; claims are the
clusters because episodes within a claim share its field). Attenuation-form comparison
uses five-fold cross-validation with folds split by claim, alongside AIC; the
form parameter ($\gamma$ for divisive normalization, the rate for load decay)
is profiled on the full data ($\gamma = 0.9$) and re-profiled inside each
training fold, and a held-out claim's fixed effect, which cannot be estimated
from the training folds, is set to zero, so the cross-validated likelihood
scores prediction for an unseen claim with an unknown field. The divisive
$\astar$ of Section~III of the main article uses $\gamma$ profiled separately
per variant (0.53 for A, 1.32 for B); Fig.~3 of the main article uses the pooled
$\gamma = 0.9$. Per-$k$ coefficients are identified by a joint logistic fit
over all inbox sizes with claim fixed effects anchored by the empty-inbox
cells: at a single $k$ the columns $1$, $l$, and $k - l$ are collinear, so
separate per-$k$ fits identify only $c_0 + k\beta_F$ and $\beta_T - \beta_F$.
Per-claim response functions for the collective claims use the reduced form
$c_0 + g(k)[\delta k + \bar\beta(2l - k)]$ ($+\theta s$ for variant B), with
a weak ridge penalty ($10^{-3}$) where $k$ takes three values. For Fig.~3 of the main article, one logistic fit per prompt variant
with claim fixed effects and message weights $\beta_T g(k)$, $\beta_F g(k)$
(divisive, $\gamma = 0.9$ fixed, or $g \equiv 1$) assigns each query a
predicted probability; a design cell's net input $u$ is the logit of its mean
predicted probability, and its observed correct-side rate is plotted against
$u$ with a Wilson 95\% interval. Fixed points of the mean-field map are
bracketed on an endpoint-inclusive grid of 4{,}001 points in $x$ with
bisection refinement and classified by $|F'(x)| < 1$; a phase is bistable
when an unstable interior point separates two stable ones, and $\astar$ is
located by bisection in $\alpha$ on that classification, with the fold of $G$
solved jointly from $G = x$ and $G_x = 1$. Prediction-side intervals resample
each claim's single-shot outcomes within that claim and within each design
cell (200 resamples; a within-claim bootstrap, not the claim-clustered
bootstrap used for pooled coefficients) and recompute the fixed points. Reanalysis statistics on the public dataset use
question-clustered bootstrap and CR0 sandwich errors.

\emph{Outcome and threshold.} An episode's outcome is its correct-side
fraction after eight rounds, $x_8$; $q(x_0)$ is the fraction of episodes with
$x_8 \ge 0.6$, a finite-horizon quantity that we compare with the
infinite-time fixed points only through surrogate rollouts. Observed
thresholds are the 50\% crossings of a logistic fit of the correct outcome on
$x_0$ with a weak ridge penalty (0.002) so that separated cells have a
finite estimate; the slope is not constrained, a fit with a non-positive slope
would be discarded, and none occurred in the reported cells. Intervals are
episode-bootstrap percentiles (500 replicates); replicates without a finite
crossing in $(-1, 2)$ are dropped, which retained all replicates in 12 of the
17 cells observed bistable, 97--99\% in three, and 89\% and 65\% in the two
cells whose crossing lies nearest the lowest seeded fraction (claims 656 and
1146 at $\alpha = 0.55$). Rank correlations between predicted and observed thresholds are
tested by permuting claim labels, claims being the exchangeable unit.

\subsection{Pre-registration timeline: full statement}

Throughout this paper, ``pre-registered'' means that a prediction or a
decision rule, together with the criteria for judging it, was written to a
file before the corresponding data were collected, and that the file's
timestamp and SHA-256 hash are recorded in the project repository, so that
the order of prediction and data is documented; these are local records, not
a third-party time stamp. No external registry (such as OSF) was used. Three items were
pre-registered: (i)~the $\astar = 0.435 \pm 0.03$ prediction and the decision
criteria of the main collective experiment, before any collective episode
(2026-08-23); (ii)~the three-hypothesis prediction table of the polarity
discrimination experiment, before its data collection (2026-08-25);
(iii)~the claim-selection rule of the eight-claim calibration campaign,
written into the campaign instructions before its data were collected
(2026-08-27). The single-agent measurements preceded the collective
experiments. The claim-resolved predictions of
Sections~IV-B to IV-D of the main article were computed after the
collective experiments, with the estimation procedure described under
Statistics, using only the single-agent data for parameter fitting; they were
not prospectively registered. The one adaptive element (second-pass allocation of episodes)
followed rules fixed in the pre-registered design.


\bibliographystyle{IEEEtran}
\bibliography{refs}
\end{document}
